% !TEX root = ../jalase02.tex
%% ============================================================
%%  جلسهٔ ۲ — فضای احتمال: قانون احتمال، σ-میدان و میدان بورل
%%  (نتایج اصول کولموگروف تا پیوستگی احتمال به جلسهٔ ۳ منتقل شده:
%%   sessions/session03-kolmogorov-results.tex)
%% ============================================================
\chapter{فضای احتمال و اصول موضوعهٔ کولموگروف}
\label{ch:s2}

\sessioncard{دوشنبه ۶ مهر}
{فضای احتمال و اصول سه‌گانهٔ کولموگروف}
{\lr{Pishro-Nik} بخش \lr{1.3} \,؛\, \lr{Ross} فصل ۲ \,؛\, \lr{Papoulis} فصل ۲}
{\begin{itemize}
  \item جایگاه احتمال را در میان سه نگاه به مجموعه‌ها (توپولوژی، جبر، نظریهٔ اندازه) توضیح دهید؛
  \item سه قانون احتمال را بیان کنید و آن‌ها را روی مثال‌های سکه، سرور و مقاومت پیاده کنید؛
  \item دیدگاه فراوانی و دیدگاه بیزی دربارهٔ معنای «احتمال» را مقایسه کنید؛
  \item توضیح دهید چرا قانون سوم باید برای اجتماع \emph{شمارا} برقرار باشد و چرا وزن اعداد گویا صفر است؛
  \item توضیح دهید چرا همیشه نمی‌توان روی همهٔ زیرمجموعه‌ها احتمال تعریف کرد، و مفهوم $\sigma$-میدان و میدان بورل را تعریف کنید.
\end{itemize}}

%% ------------------------------------------------------------
\section{مرور: برآمدها و فضای نمونه}
%% ------------------------------------------------------------

در جلسهٔ اول گفتیم هر نتیجهٔ ممکن یک آزمایش تصادفی را یک \term{برآمد}{outcome} می‌نامیم و همهٔ برآمدها را در یک مجموعه به نام \term{فضای نمونه}{sample space} جمع می‌کنیم و آن را $\Omega$ می‌نامیم. هر زیرمجموعهٔ $A\subseteq\Omega$ یک \term{پیشامد}{event} است و «$A$ رخ داده» یعنی برآمدِ آزمایش داخل $A$ افتاده است.

فضای نمونه از نظر «تعداد اعضا» سه حالت دارد و، همان‌طور که خواهیم دید، رفتار احتمال در این سه حالت یکسان نیست:

\begin{figblock}
\begin{tikzpicture}
  % متناهی
  \begin{scope}[xshift=11cm]
    \draw[defcolor, thick, fill=defcolor!7, rounded corners=8pt] (-2.4,-1.5) rectangle (2.4,1.5);
    \node[font=\small] at (0,0.35) {\lr{$\{H,\,T\}$}};
    \node[font=\small] at (0,-0.45) {\lr{$\{1,2,\dots,6\}$}};
    \node[font=\headfont, text=defcolor] at (0,2.05) {\rl{متناهی}};
    \node[font=\scriptsize, text=graytext] at (0,-2.0) {\rl{سکه، تاس}};
  \end{scope}
  % شمارا
  \begin{scope}[xshift=5.5cm]
    \draw[excolor, thick, fill=excolor!7, rounded corners=8pt] (-2.4,-1.5) rectangle (2.4,1.5);
    \foreach \i in {0,...,4} \fill[excolor] ({-1.8+0.7*\i},0) circle (2.3pt);
    \node[font=\small] at (1.8,0) {\lr{$\cdots$}};
    \node[font=\scriptsize] at (-1.8,-0.45) {\lr{$0$}}; \node[font=\scriptsize] at (-1.1,-0.45) {\lr{$1$}}; \node[font=\scriptsize] at (-0.4,-0.45) {\lr{$2$}};
    \node[font=\headfont, text=excolor] at (0,2.05) {\rl{نامتناهیِ شمارا}};
    \node[font=\scriptsize, text=graytext] at (0,-2.0) {\rl{تعداد درخواست‌های یک سرور}};
  \end{scope}
  % ناشمارا
  \begin{scope}[xshift=0cm]
    \draw[thmcolor, thick, fill=thmcolor!7, rounded corners=8pt] (-2.4,-1.5) rectangle (2.4,1.5);
    \draw[thmcolor, line width=2.5pt] (-1.7,0) -- (1.7,0);
    \node[font=\scriptsize] at (-1.7,-0.45) {\lr{$0.5$}}; \node[font=\scriptsize] at (1.7,-0.45) {\lr{$1.5$}};
    \node[font=\headfont, text=thmcolor] at (0,2.05) {\rl{ناشمارا}};
    \node[font=\scriptsize, text=graytext] at (0,-2.0) {\rl{مقدار یک مقاومت، زمان، ولتاژ}};
  \end{scope}
\end{tikzpicture}

\figcaption[fig:omega-types]{سه نوع فضای نمونه. در این جلسه برای هر کدام یک مثال کامل می‌بینیم.}
\end{figblock}

%% ------------------------------------------------------------
\section{سه نگاه به یک مجموعه: توپولوژی، جبر، نظریهٔ اندازه}
%% ------------------------------------------------------------

وقتی یک مجموعه در دست داریم، سه کار کاملاً متفاوت می‌توانیم با آن انجام دهیم؛ و هر کدام یک شاخهٔ ریاضی است:

\begin{figblock}
\begin{tikzpicture}[
  sci/.style={draw=#1, line width=1.2pt, fill=#1!7, rounded corners=8pt, text width=4.4cm, align=center,
              minimum height=2.9cm, font=\small},
  arr/.style={-{Stealth[length=3mm]}, line width=1.1pt, graytext}]
  \node[circle, draw=primary, line width=1.5pt, fill=primarylight, minimum size=2.1cm, font=\headfont\Large] (O) at (0,0) {\lr{$\Omega$}};
  \node[sci=defcolor] (T) at (5.4,1.9) {\rl{{\headfont\color{defcolor} توپولوژی (هندسه)}}\\[2pt] \rl{فاصله و نزدیکی نقاط، شکل‌ها، مجموعه‌های باز}\\[2pt] \lr{$d(x,y)=|x-y|$}};
  \node[sci=excolor] (A) at (-5.4,1.9) {\rl{{\headfont\color{excolor} جبر}}\\[2pt] \rl{عملیات روی مجموعه‌ها}\\[2pt] \lr{$A\cup B,\ A\cap B,\ A^c$}};
  \node[sci=thmcolor, line width=2.4pt] (M) at (0,-3.5) {\rl{{\headfont\color{thmcolor} نظریهٔ اندازه}}\\[2pt] \rl{وزن و ارزش دادن به مجموعه‌ها}\\[2pt] \rl{طول، مساحت، شمارش، \textbf{احتمال}}};
  \draw[arr] (O) -- (T); \draw[arr] (O) -- (A); \draw[arr, thmcolor, line width=1.8pt] (O) -- (M);
\end{tikzpicture}

\figcaption[fig:three-sciences]{سه نگاه به یک مجموعه. نظریهٔ احتمال در دستهٔ سوم است و از جبر مجموعه‌ها به عنوان ابزار استفاده می‌کند.}
\end{figblock}

\begin{enumerate}
  \item \textbf{توپولوژی (هندسه):} به کمک \emph{فاصله} بین نقاط، مفاهیمی مانند نزدیک بودن، شکل، مجموعهٔ باز و پیوستگی را تعریف می‌کند.
  \item \textbf{جبر:} روی مجموعه‌ها \emph{عملیات} انجام می‌دهد: اجتماع، اشتراک، متمم، تفاضل. (همان چیزی که در جلسهٔ ۱ مرور کردیم.)
  \item \textbf{نظریهٔ اندازه:} به مجموعه‌ها \emph{وزن} یا \emph{ارزش} نسبت می‌دهد؛ مثل طول یک بازه، مساحت یک ناحیه یا تعداد اعضای یک مجموعه.
\end{enumerate}

\begin{keyformula}
\[ \text{\rl{در احتمال، کار سوم را انجام می‌دهیم: به پیشامدها (زیرمجموعه‌های $\Omega$) وزن می‌دهیم.}} \]
\end{keyformula}

احتمال پیشامد $A$ را با $\P(A)$ نشان می‌دهیم. تصویر ذهنی خوبی که در سراسر این جلسه به کار می‌بریم: \key{یک کیلو شن را روی $\Omega$ پخش کرده‌ایم؛ $\P(A)$ مقدار شنی است که روی ناحیهٔ $A$ نشسته است.}

\subsection{وزن‌هایی که از قبل می‌شناسیم}

\begin{center}
\renewcommand{\arraystretch}{1.35}
\begin{tabular}{>{\semifont}l l l}
\toprule
نوع وزن & به چه مجموعه‌هایی & چه عددی می‌دهد \\
\midrule
شمارش & زیرمجموعه‌های یک مجموعهٔ متناهی & تعداد اعضا: $|A|$\\
طول & بازه‌های $\R$ & $b-a$ برای بازهٔ $[a,b]$\\
مساحت & شکل‌های صفحه & مثلاً طول $\times$ عرض برای مستطیل\\
جرم & ناحیه‌های یک جسم & جرم آن ناحیه\\
\bottomrule
\end{tabular}
\end{center}

همهٔ این وزن‌ها سه ویژگی مشترک دارند: \emph{منفی نمی‌شوند}، \emph{به کل فضا مقدار مشخصی می‌دهند}، و \emph{روی تکه‌های جدا از هم جمع می‌شوند} (طول $[0,3]$ برابر طول $[0,1]$ به علاوهٔ طول $(1,3]$ است). هر وزنی با این ویژگی‌ها را یک \term{اندازه}{measure} می‌نامیم. احتمال هم یک اندازه است؛ اندازه‌ای که وزن کل فضا در آن ۱ است.

%% ------------------------------------------------------------
\section{قانون احتمال}
%% ------------------------------------------------------------

از یک «وزن‌دهی» معقول به پیشامدها چه انتظاری داریم؟ با تصویر شن:
(۱) مقدار شن روی هیچ ناحیه‌ای منفی نیست؛ (۲) کل شن یک کیلوست؛ (۳) اگر دو ناحیه هم‌پوشانی نداشته باشند، شنِ روی اجتماعشان برابر جمع شنِ روی هر کدام است. همین سه انتظار را قانون می‌کنیم:

\begin{definition}[label=def:law-v1]{قانون احتمال (نسخهٔ اول)}
به هر پیشامد $A$ عدد $\P(A)$ نسبت می‌دهیم به طوری که:
\begin{enumerate}[leftmargin=5.4em]
  \item[\textbf{قانون ۱.}] \textbf{وزن منفی نیست:} \quad $\P(A)\ge0$؛
  \item[\textbf{قانون ۲.}] \textbf{وزن کل فضا یک است:} \quad $\P(\Omega)=1$؛
  \item[\textbf{قانون ۳.}] \textbf{جمع‌پذیری:} \quad اگر $A$ و $B$ مجزا باشند ($A\cap B=\varnothing$)، آنگاه $\P(A\cup B)=\P(A)+\P(B)$.
\end{enumerate}
این سه قانون را \term{قانون احتمال}{law of probability} می‌نامیم.
\end{definition}

کل نظریهٔ احتمال روی همین سه قانون سوار می‌شود: هر قاعده‌ای که بعداً می‌بینیم (قاعدهٔ متمم، شمول و عدم شمول، قضیهٔ بیز و \dots) از همین‌ها نتیجه می‌شود. حالا با سه مثال این قانون‌ها را جا می‌اندازیم؛ و در مسیر خواهیم دید که قانون ۳ باید کمی قوی‌تر شود.

\begin{history}{کولموگروف، ۱۹۳۳}
تا اوایل قرن بیستم احتمال تعریف ریاضی دقیقی نداشت و هر کس برداشت خودش را داشت (کلاسیک، فراوانی، ذهنی). آندری کولموگروف در کتاب «مبانی نظریهٔ احتمال» (۱۹۳۳) پیشنهاد کرد به جای بحث فلسفی دربارهٔ «معنای» احتمال، فقط \emph{قواعد} آن را مشخص کنیم؛ و این قواعد را از نظریهٔ اندازه (بورل و لبگ) گرفت. این قواعد امروز «اصول موضوعهٔ کولموگروف» نام دارند.
\end{history}

%% ------------------------------------------------------------
\section{مثال ۱: پرتاب سکه}
%% ------------------------------------------------------------

\begin{example}[label=ex:coin-law]{قانون احتمال برای یک سکه}
سکه‌ای (نه لزوماً سالم) را یک بار پرتاب می‌کنیم. فضای نمونه $\Omega=\{H,T\}$ است (شیر $=H$، خط $=T$) و چهار پیشامد داریم: $\varnothing$، $\{H\}$، $\{T\}$ و $\Omega$. قانون احتمال را روی آن پیاده کنید.
\begin{solution}
فرض کنید احتمال شیر را $p$ بنامیم: $\P(\{H\})=p$. بقیه از قانون‌ها به دست می‌آیند:
\begin{itemize}
  \item $\{H\}$ و $\{T\}$ مجزا هستند و اجتماعشان $\Omega$ است. طبق قانون‌های ۳ و ۲: \ $\P(\{H\})+\P(\{T\})=\P(\Omega)=1$؛ پس $\P(\{T\})=1-p$.
  \item $\Omega$ و $\varnothing$ مجزا هستند و $\Omega\cup\varnothing=\Omega$. طبق قانون ۳: \ $\P(\Omega)=\P(\Omega)+\P(\varnothing)$؛ پس $\P(\varnothing)=0$.
  \item طبق قانون ۱، $p\ge0$ و $1-p\ge0$؛ پس $0\le p\le1$.
\end{itemize}
\begin{center}
\renewcommand{\arraystretch}{1.3}
\begin{tabular}{c|cccc}
پیشامد & $\varnothing$ & $\{H\}$ & $\{T\}$ & $\Omega$\\\hline
احتمال & $0$ & $p$ & $1-p$ & $1$
\end{tabular}
\end{center}
برای هر $p\in[0,1]$ این جدول هر سه قانون را برآورده می‌کند؛ یعنی \emph{یک خانوادهٔ کامل} از مدل‌های معتبر داریم.
\end{solution}
\end{example}

\begin{note}
قانون‌های احتمال به ما نمی‌گویند $p$ \emph{چند} است؛ فقط می‌گویند وزن‌ها باید با هم سازگار باشند. انتخاب $p$ کار \emph{مدل‌سازی} یا \emph{آمار} است. اما خودِ عدد $p$ چه معنایی دارد؟ دو دیدگاه مشهور وجود دارد.
\end{note}

\subsection{دیدگاه فراوانی: احتمال یعنی نسبت تکرار}

سکه را $N$ بار (تحت شرایط یکسان) پرتاب می‌کنیم و $N_H$ را تعداد شیرها می‌گیریم. در این دیدگاه:
\[ p=\lim_{N\to\infty}\frac{N_H}{N}. \]
یعنی «احتمال شیر ۰٫۵ است» به این معناست که اگر سکه را بارها و بارها بیندازیم، تقریباً نیمی از نتایج شیر است. (آزمایش‌های تاریخی بوفون با ۴۰۴۰ پرتاب و پیرسون با ۲۴۰۰۰ پرتاب را در جلسهٔ ۱ دیدیم.)

\textbf{محدودیت:} این دیدگاه فقط برای آزمایش‌های \emph{تکرارپذیر} معنا دارد. «احتمال اینکه این پروژه تا پایان ترم تمام شود» را نمی‌توان با هزار بار تکرار سنجید.

\subsection{دیدگاه بیزی: احتمال یعنی میزان باور}

در این دیدگاه، $p$ \emph{میزان باور} ما به شیر آمدن است؛ در اصل اعتقاد ما این است البته  این اعتقاد با دیدن داده‌های تازه \emph{به‌روز} می‌شود. مثلاً پیش از هر پرتاب، به یک سکهٔ معمولی باور داریم که $p$ حدود $0.5$ است. اگر در ۱۰ پرتاب ۸ بار شیر ببینیم، منطقی است باورمان به سمت $p$های بزرگ‌تر جابه‌جا شود. ابزار دقیق این به‌روزرسانی \emph{قضیهٔ بیز} است (جلسهٔ ۳) و استنتاج بیزی را در جلسه‌های ۲۶ و ۲۷ به تفصیل می‌بینیم.

\begin{intuition}
دو دیدگاه دربارهٔ \emph{معنای} احتمال اختلاف دارند، ولی هر دو از \emph{یک قانون} پیروی می‌کنند: نسبت‌های تکرار $N_H/N$ منفی نیستند، کلشان ۱ است و برای پیشامدهای مجزا جمع می‌شوند؛ باورهای معقول هم همین‌طور. به همین دلیل همهٔ قضیه‌هایی که از قانون احتمال ثابت می‌کنیم، در هر دو دیدگاه معتبرند.
\end{intuition}

%% ------------------------------------------------------------
\section{مثال ۲: درخواست‌های یک سرور و ارتقای قانون سوم}
%% ------------------------------------------------------------

\begin{example}[label=ex:server]{تعداد درخواست‌ها در یک ثانیه}
می‌خواهیم تعداد درخواست‌هایی را که در یک ثانیه به یک سرور می‌رسد مدل کنیم. فضای نمونه $\Omega=\{0,1,2,\dots\}$ نامتناهیِ شماراست. فرض کنید مدل ما این باشد:
\[ \P(\{i\})=\frac{1}{2^{\,i+1}},\qquad i=0,1,2,\dots \]
یعنی $\P(\{0\})=\tfrac12$، \ $\P(\{1\})=\tfrac14$، \ $\P(\{2\})=\tfrac18$ و الی آخر. احتمال اینکه تعداد درخواست‌ها \emph{زوج} باشد چقدر است؟
\end{example}

\begin{figblock}
\begin{tikzpicture}
\begin{axis}[
  width=10.5cm, height=4.4cm, ybar, bar width=14pt,
  ymin=0, ymax=0.6, xtick={0,...,7}, ytick=\empty, axis y line=none, axis x line*=bottom,
  tick label style={font=\footnotesize},
  nodes near coords, nodes near coords style={font=\scriptsize, color=graytext},
  point meta=explicit symbolic]
  \addplot[fill=excolor!55, draw=excolor] coordinates {
    (0,0.5)[$\frac12$] (1,0.25)[$\frac14$] (2,0.125)[$\frac18$] (3,0.0625)[$\frac1{16}$]
    (4,0.03125)[$\frac1{32}$] (5,0.015625)[] (6,0.0078125)[] (7,0.00390625)[]};
\end{axis}
\end{tikzpicture}

\figcaption[fig:server]{وزن هر تعداد درخواست. برای هر $i$ وزنی داریم؛ پس \emph{بی‌نهایت} تکه وزن داریم.}
\end{figblock}

\subsection{قانون سوم با «دو» پیشامد کافی نیست}

قانون ۳ (نسخهٔ اول) برای \emph{دو} پیشامد مجزاست. با استقرا می‌توان آن را به هر تعداد \emph{متناهی} پیشامد مجزا تعمیم داد:
\[ \P(A_1\cup\cdots\cup A_n)=\P(A_1)+\cdots+\P(A_n). \]
اما پیشامد «تعداد درخواست‌ها زوج است» برابر $\{0,2,4,6,\dots\}=\{0\}\cup\{2\}\cup\{4\}\cup\cdots$ است: اجتماعی از \emph{بی‌نهایت} پیشامد مجزا. متمم آن (تعداد فرد) هم همین مشکل را دارد. پس با نسخهٔ اول قانون‌ها \emph{هیچ راهی} برای محاسبهٔ احتمال آن نداریم!

حتی بدتر: خودِ اینکه وزن‌های $\tfrac12,\tfrac14,\tfrac18,\dots$ یک مدل معتبرند (یعنی جمعشان با $\P(\Omega)=1$ سازگار است) به همین مشکل برمی‌خورد، چون $\Omega=\{0\}\cup\{1\}\cup\{2\}\cup\cdots$ هم اجتماعی نامتناهی است. طبیعی است که بخواهیم:
\[ \P(A_1\cup A_2\cup\cdots)=\sum_{i=1}^\infty\P(A_i)\qquad\text{\rl{برای $A_i$های دوبه‌دو مجزا.}} \]

\begin{definition}[label=def:law-final]{قانون احتمال (نسخهٔ نهایی)}
قانون‌های ۱ و ۲ مثل قبل‌اند و قانون ۳ به این صورت تقویت می‌شود:
\begin{enumerate}[leftmargin=5.4em]
  \item[\textbf{قانون ۳.}] \textbf{جمع‌پذیری شمارا:} اگر $A_1,A_2,\dots$ پیشامدهایی \emph{دوبه‌دو مجزا} باشند ($A_i\cap A_j=\varnothing$ برای $i\ne j$)، آنگاه
  \[ \P\Bigl(\bigcup_{i=1}^\infty A_i\Bigr)=\sum_{i=1}^\infty\P(A_i). \]
\end{enumerate}
\end{definition}

\begin{keyformula}
\[ \P(A)\ge0,\qquad \P(\Omega)=1,\qquad A_i\cap A_j=\varnothing\ (i\ne j)\ \Rightarrow\ \P\Bigl(\bigcup_i A_i\Bigr)=\sum_i\P(A_i) \]
\end{keyformula}

\begin{note}
نسخهٔ نهایی، نسخهٔ اول را هم در بر دارد: کافی است $A_1=A$، $A_2=B$ و $A_3=A_4=\cdots=\varnothing$ بگیریم (و از $\P(\varnothing)=0$ استفاده کنیم که مثل مثال سکه اثبات می‌شود).
\end{note}

\begin{solution}
(ادامهٔ مثال \ref{ex:server}) حالا مسئله حل می‌شود. طبق جمع‌پذیری شمارا:
\[ \P(\Omega)=\sum_{i=0}^\infty\frac1{2^{\,i+1}}=\frac12+\frac14+\frac18+\cdots=1\ \checkmark \]
\[ \P(\fa{تعداد زوج})=\sum_{k=0}^\infty\P(\{2k\})=\frac12+\frac18+\frac1{32}+\cdots=\frac{1/2}{1-1/4}=\frac23. \]
به همین ترتیب احتمال «بیش از ۳ درخواست» برابر $\sum_{i=4}^\infty 2^{-(i+1)}=\tfrac1{16}$ است.
\end{solution}

%% ------------------------------------------------------------
\section{مثال ۳: ساختن یک مقاومت ۱ اُهمی}
%% ------------------------------------------------------------

\begin{example}[label=ex:resistor]{مقاومت با تلرانس}
کارخانه‌ای مقاومت‌های «۱ اُهمی» می‌سازد، ولی به دلیل خطای ساخت، مقدار واقعی هر مقاومت عددی در بازهٔ $[0.5,\,1.5]$ اُهم است و هیچ مقداری از این بازه بر دیگری ترجیح ندارد. قانون احتمال را چگونه روی این آزمایش پیاده کنیم؟
\end{example}

فضای نمونه $\Omega=[0.5,\,1.5]$ است؛ یک مجموعهٔ \emph{ناشمارا}. اینجا دیگر نمی‌توانیم مثل مثال ۲ به هر نقطه وزن بدهیم و جمع بزنیم. طبیعی‌ترین مدل این است: \key{احتمال هر بازه برابر طول آن است.}
\[ \P\bigl([a,b]\bigr)=b-a,\qquad 0.5\le a\le b\le1.5. \]
(طول کل فضا ۱ است، پس قانون ۲ برقرار است.) مثلاً احتمال اینکه مقاومت بین $0.9$ و $1.1$ اُهم باشد $0.2$ است.

\begin{figblock}
\begin{tikzpicture}[x=10cm, y=1cm]
  \fill[accent!18] (0.5,0) rectangle (1.5,0.6);
  \fill[accent!70] (0.9,0) rectangle (1.1,0.6);
  \draw[primary, thick] (0.5,0) rectangle (1.5,0.6);
  \foreach \t in {0.5,0.9,1,1.1,1.5} { \draw[primary] (\t,0) -- (\t,-0.12); \node[font=\small, below] at (\t,-0.12) {\lr{$\t$}}; }
  \fill[warncolor] (1,0.3) circle (2.2pt);
  \draw[decorate, decoration={brace, amplitude=5pt}, primary] (0.9,0.7) -- (1.1,0.7) node[midway, above=5pt, font=\small] {\lr{$\mathbb P=0.2$}};
  \node[font=\small, warncolor] at (1.36,1.12) {\rl{نقطهٔ ۱: احتمال صفر}};
  \draw[warncolor, -{Stealth}] (1.24,0.98) -- (1.015,0.36);
\end{tikzpicture}

\figcaption[fig:resistor]{مدل مقاومت: شن به طور یکنواخت روی $[0.5,1.5]$ پخش شده است. بازه‌ها وزن مثبت دارند ولی هر نقطهٔ منفرد وزن صفر.}
\end{figblock}

\subsection{وزن یک نقطه: چرا «دقیقاً ۱ اُهم» احتمال صفر دارد؟}

پیشامد «مقاومت \emph{دقیقاً} ۱ اُهم است» یعنی $\{1\}$. برای هر $\varepsilon>0$ داریم $\{1\}\subseteq[1-\varepsilon,\,1+\varepsilon]$. اگر $B\subseteq A$، آنگاه $A=B\cup(A-B)$ و دو تکه مجزا هستند؛ پس طبق قانون‌های ۳ و ۱، $\P(A)=\P(B)+\P(A-B)\ge\P(B)$. بنابراین
\[ 0\le\P(\{1\})\le\P\bigl([1-\varepsilon,1+\varepsilon]\bigr)=2\varepsilon\quad\text{\rl{برای هر $\varepsilon>0$}}\qquad\Longrightarrow\qquad \P(\{1\})=0. \]
همین استدلال برای هر نقطهٔ $x\in[0.5,1.5]$ کار می‌کند: \ $\P(\{x\})=0$.

\begin{intuition}
در عمل هیچ کارخانه‌ای مقاومتِ «دقیقاً» ۱ اُهمی نمی‌سازد؛ همیشه با خطای ساخت سروکار داریم. به همین دلیل در مشخصات قطعات، \emph{تلرانس} (یک بازه) نوشته می‌شود و نه یک عدد دقیق. «احتمال صفر» برای یک نقطه یعنی همین.
\end{intuition}

\subsection{وزن اعداد گویا صفر است و وزن اعداد گنگ یک}

حالا می‌پرسیم: احتمال اینکه مقدار مقاومت یک عدد \emph{گویا} باشد چقدر است؟ اعداد گویای $[0.5,1.5]$ \emph{شمارا} هستند (جلسهٔ ۱)؛ پس می‌توانیم آن‌ها را در یک دنباله بنویسیم: $q_1,q_2,q_3,\dots$. پیشامد «گویا بودن» اجتماعی شمارا از پیشامدهای مجزای $\{q_i\}$ است و طبق جمع‌پذیری شمارا:
\[ \P(\fa{گویا})=\P\Bigl(\bigcup_{i=1}^\infty\{q_i\}\Bigr)=\sum_{i=1}^\infty\P(\{q_i\})=\sum_{i=1}^\infty0=0. \]
پیشامد «گنگ بودن» متمم آن است. چون گویاها و گنگ‌ها مجزا هستند و اجتماعشان $\Omega$ است، طبق قانون‌های ۲ و ۳:
\[ \P(\fa{گنگ})=1-\P(\fa{گویا})=1-0=1. \]

\begin{warning}
شاید به نظر برسد تناقضی وجود دارد: هر نقطه احتمال صفر دارد، ولی بازهٔ $[0.9,1.1]$ که از نقطه‌ها ساخته شده احتمال $0.2$ دارد! تناقضی نیست، چون یک بازه اجتماع \emph{ناشمارا} نقطه است و قانون ۳ فقط برای اجتماع \emph{شمارا} برقرار است. اگر قانون ۳ را برای اجتماع ناشمارا هم می‌خواستیم، نتیجه می‌گرفتیم $1=\P(\Omega)=\sum_{x\in\Omega}\P(\{x\})=0$. پس «شمارا» دقیقاً انتخاب درست است: متناهی برای مثال سرور کم است و ناشمارا به تناقض می‌رسد.
\end{warning}

\subsection{آیا می‌توانیم به «همهٔ» زیرمجموعه‌ها وزن بدهیم؟}

در مثال‌های ۱ و ۲ به \emph{همهٔ} زیرمجموعه‌های $\Omega$، یعنی مجموعهٔ توانی $2^\Omega$، احتمال دادیم. در مثال ۳ هم می‌خواهیم «وزن = طول» را برای \emph{هر} زیرمجموعهٔ $[0.5,1.5]$ تعریف کنیم؛ نه فقط برای بازه‌ها، نقطه‌ها و گویاها. شگفت آنکه این کار \emph{ممکن نیست}:

\begin{moreinfo}{قضیهٔ ویتالی (۱۹۰۵)}
\textbf{صورت قضیه.} هیچ وزن‌دهی $\lambda$ روی \emph{همهٔ} زیرمجموعه‌های $[0,1)$ وجود ندارد که هم‌زمان (۱) جمع‌پذیری شمارا داشته باشد، (۲) به هر بازه طولش را بدهد، و (۳) با «چرخاندن» مجموعه (جابه‌جا کردن به پیمانهٔ ۱، مثل چرخاندن روی دایره‌ای به محیط ۱) تغییر نکند. (برای $[0.5,1.5]$ هم دقیقاً همین‌طور است؛ فقط بازه را نیم واحد جابه‌جا کنید.)

\textbf{ایدهٔ اثبات.}
\begin{enumerate}
  \item دو عدد $x,y\in[0,1)$ را «هم‌خانواده» بگیرید اگر $x-y$ گویا باشد. بازه به خانواده‌های مجزا افراز می‌شود.
  \item از \emph{هر} خانواده دقیقاً یک نماینده بردارید و مجموعهٔ نماینده‌ها را $V$ بنامید (این انتخاب با «اصل انتخاب» ممکن است).
  \item برای هر گویای $q\in[0,1)$، کپی چرخیدهٔ $V_q=V\oplus q$ را بسازید. این کپی‌ها \emph{دوبه‌دو مجزا} هستند، تعدادشان \emph{شماراست} و اجتماعشان \emph{کل} $[0,1)$ است.
  \item همهٔ کپی‌ها وزن برابری مثل $w$ دارند (فقط چرخیده‌اند). طبق جمع‌پذیری شمارا: $1=w+w+w+\cdots$. اگر $w=0$ جمع صفر است و اگر $w>0$ جمع بی‌نهایت. \textbf{تناقض.}
\end{enumerate}

\textbf{نتیجه.} مجموعهٔ $V$ (\emph{مجموعهٔ ویتالی}) آن‌قدر «بدقواره» است که نمی‌توان به آن وزن داد. پس روی فضای ناشمارا \emph{نمی‌توانیم همیشه} قانون احتمال را روی کل $2^\Omega$ تعریف کنیم.
(از همین خانواده است پارادوکس مشهور باناخ--تارسکی (۱۹۲۴): یک گوی را می‌توان به ۵ تکه تقسیم کرد و با جابه‌جا کردن آن‌ها دو گوی کامل ساخت!)
\end{moreinfo}

%% ------------------------------------------------------------
\section{عقب‌نشینی: $\sigma$-میدان}
%% ------------------------------------------------------------

قضیهٔ ویتالی ما را مجبور به \emph{عقب‌نشینی} می‌کند: بعضی از زیرمجموعه‌های $\Omega$ هستند که نمی‌توان به آن‌ها وزن داد. پس به جای همهٔ زیرمجموعه‌ها، یک خانوادهٔ $\F$ از زیرمجموعه‌های $\Omega$ انتخاب می‌کنیم؛ \key{زیرمجموعه‌هایی که بتوان قانون احتمال را روی آن‌ها سوار کرد.} اعضای $\F$ را \emph{پیشامد} می‌نامیم.

\subsection{چه ویژگی‌هایی لازم داریم؟}

این خانواده نمی‌تواند دلخواه باشد؛ باید با قانون احتمال و منطق معمولی ما جور باشد:

\begin{center}
\renewcommand{\arraystretch}{1.45}
\begin{tabular}{l l}
\toprule
\textbf{چرا لازم است} & \textbf{ویژگی $\F$}\\
\midrule
قانون ۲ دربارهٔ $\P(\Omega)$ حرف می‌زند؛ پس $\Omega$ باید پیشامد باشد & $\Omega\in\F$\\
اگر بتوان پرسید «$A$ رخ داد؟»، باید بتوان پرسید «$A$ رخ نداد؟» & $A\in\F\Rightarrow A^c\in\F$\\
قانون ۳ دربارهٔ $\P\bigl(\bigcup_iA_i\bigr)$ حرف می‌زند؛ پس اجتماع باید پیشامد باشد & $A_1,A_2,\ldots\in\F\Rightarrow\bigcup_iA_i\in\F$\\
\bottomrule
\end{tabular}
\end{center}

\begin{definition}[label=def:sigma]{$\sigma$-میدان ($\sigma$-جبر)}
خانوادهٔ $\F$ از زیرمجموعه‌های $\Omega$ یک \term{$\sigma$-میدان}{$\sigma$-field} (یا \term{$\sigma$-جبر}{$\sigma$-algebra}) است اگر:
\begin{enumerate}
  \item $\Omega\in\F$؛
  \item \textbf{بسته بودن نسبت به متمم:} اگر $A\in\F$ آنگاه $A^c\in\F$؛
  \item \textbf{بسته بودن نسبت به اجتماع شمارا:} اگر $A_1,A_2,\dots\in\F$ آنگاه $\bigcup_{i=1}^\infty A_i\in\F$.
\end{enumerate}
اعضای $\F$ را \emph{پیشامد} می‌نامیم.
\end{definition}

حرف $\sigma$ (سیگما) در این نام یادآور «شمارا»ست؛ همان‌طور که $\sum$ نماد جمع است، $\sigma$ به \emph{اجتماع شمارا} اشاره می‌کند.

\begin{proposition}[label=prop:sigma-props]{خواص $\sigma$-میدان}
اگر $\F$ یک $\sigma$-میدان باشد، آنگاه:
(الف) $\varnothing\in\F$؛
(ب) $\F$ نسبت به اجتماع \emph{متناهی} بسته است؛
(ج) $\F$ نسبت به اشتراک شمارا (و متناهی) بسته است؛
(د) اگر $A,B\in\F$ آنگاه $A-B\in\F$.
\end{proposition}

\begin{proof}
(الف) $\varnothing=\Omega^c$. \
(ب) برای $A_1,\dots,A_n\in\F$ قرار دهید $A_{n+1}=A_{n+2}=\cdots=\varnothing$؛ اجتماع شمارای این دنباله همان $\bigcup_{i=1}^n A_i$ است. \
(ج) طبق دمورگان $\bigcap_i A_i=\bigl(\bigcup_i A_i^c\bigr)^c$ و طرف راست با ویژگی‌های (۲) و (۳) در $\F$ است. \
(د) $A-B=A\cap B^c$.
\end{proof}

\begin{intuition}
یعنی هر کاری که با «و»، «یا»، «نه» روی تعدادی شمارا پیشامد انجام دهید، از $\F$ بیرون نمی‌روید. تعبیر دیگر: $\F$ مجموعهٔ \emph{پرسش‌هایی} است که پس از آزمایش می‌توان به آن‌ها جواب «بله/خیر» داد. مثلاً اگر تاسی بیندازند و فقط بگویند زوج آمده یا فرد، پرسش‌های قابل جواب $\F=\bigl\{\varnothing,\{1,3,5\},\{2,4,6\},\Omega\bigr\}$ است و $\{6\}$ در آن نیست.
\end{intuition}

\begin{example}[label=ex:sigma-examples]{چند $\sigma$-میدان}
\begin{enumerate}
  \item $\F=\{\varnothing,\Omega\}$: کوچک‌ترین $\sigma$-میدان.
  \item برای یک $A$ ثابت: $\F=\{\varnothing,A,A^c,\Omega\}$.
  \item $\F=2^\Omega$: بزرگ‌ترین $\sigma$-میدان. \key{در مثال‌های ۱ و ۲ (فضای متناهی یا شمارا) همین را انتخاب می‌کنیم.}
  \item روی $\Omega=\{1,2,3,4\}$، خانوادهٔ $\bigl\{\varnothing,\{1\},\{2,3,4\},\{3,4\},\{1,2\},\Omega\bigr\}$ یک $\sigma$-میدان \emph{نیست}: $\{1\}$ و $\{3,4\}$ در آن هستند ولی اجتماعشان $\{1,3,4\}$ نیست.
\end{enumerate}
\end{example}

\subsection{قانون احتمال روی $\sigma$-میدان: فضای احتمال}

حالا می‌توانیم قانون احتمال را دقیق بنویسیم: وزن را فقط به اعضای $\F$ می‌دهیم.

\begin{definition}[label=def:prob-space]{فضای احتمال}
\term{فضای احتمال}{probability space} یک سه‌تایی $(\Omega,\F,\P)$ است که در آن $\Omega$ فضای نمونه، $\F$ یک $\sigma$-میدان روی $\Omega$، و $\P:\F\to\R$ تابعی است که قانون احتمال (نسخهٔ نهایی) را برآورده می‌کند:
\[ \P(A)\ge0\ \ (A\in\F),\qquad \P(\Omega)=1,\qquad \P\Bigl(\bigcup_{i}A_i\Bigr)=\sum_{i}\P(A_i)\ \ \text{\rl{برای $A_i\in\F$ دوبه‌دو مجزا.}} \]
این سه شرط را \term{اصول موضوعهٔ کولموگروف}{Kolmogorov axioms} می‌نامیم.
\end{definition}

\begin{warning}
$\sigma$-میدان بودن \emph{شرط لازم} است، نه کافی. قانون احتمال همیشه روی یک $\sigma$-میدان تعریف می‌شود، اما \emph{هر} $\sigma$-میدانی برای \emph{هر} قانون احتمالی مناسب نیست. مثال: $2^{[0.5,1.5]}$ یک $\sigma$-میدان است، ولی طبق قضیهٔ ویتالی قانون «وزن = طول» را نمی‌توان روی آن تعریف کرد. پس $\F$ را باید طوری انتخاب کنیم که
\begin{itemize}
  \item آن‌قدر \emph{بزرگ} باشد که همهٔ پیشامدهایی که برایمان مهم‌اند (بازه‌ها، نقطه‌ها، گویاها، \dots) در آن باشند؛
  \item و آن‌قدر \emph{کوچک} باشد که قانون احتمالی که می‌خواهیم، روی آن واقعاً وجود داشته باشد.
\end{itemize}
در حالت گسسته $2^\Omega$ هر دو شرط را دارد. برای بازه‌های پیوسته جواب، میدان بورل است.
\end{warning}

%% ------------------------------------------------------------
\section{میدان بورل: $\sigma$-میدان مناسب برای بازه‌های پیوسته}
%% ------------------------------------------------------------

به مثال مقاومت برگردیم: $\Omega=[0.5,1.5]$. کدام $\sigma$-میدان را انتخاب کنیم؟ از پرسش‌هایی شروع می‌کنیم که \emph{حتماً} می‌خواهیم جواب بدهیم: «مقدار مقاومت بین $a$ و $b$ است؟». پس همهٔ بازه‌های
\[ (a,b),\qquad 0.5\le a<b\le1.5 \]
را در $\F$ می‌ریزیم. اما $\F$ باید $\sigma$-میدان باشد، پس \emph{خودبه‌خود} مجموعه‌های دیگری هم وارد آن می‌شوند.

\subsection{چه چیزهایی خودبه‌خود وارد می‌شوند؟}

\textbf{نقطه‌ها.} بازه‌های $\bigl(a-\tfrac1n,\,a+\tfrac1n\bigr)$ همه در $\F$ هستند؛ پس اشتراک شمارای آن‌ها هم باید در $\F$ باشد:
\[ \bigcap_{n=1}^\infty\Bigl(a-\frac1n,\ a+\frac1n\Bigr)=\{a\}. \]
هر عددی به جز $a$ برای $n$ به اندازهٔ کافی بزرگ از این بازه‌ها بیرون می‌افتد؛ پس فقط خود $a$ می‌ماند.

\begin{figblock}
\begin{tikzpicture}[x=1cm, y=1cm]
  \draw[primary, thick] (-6,0) -- (6,0);
  \foreach \n/\op in {1/25, 2/40, 3/55, 4/70, 6/85, 12/100} {
    \pgfmathsetmacro{\w}{5/\n}
    \draw[accent!\op, line width=4pt] ({-\w},{0.28*\n^0.8+0.25}) -- ({\w},{0.28*\n^0.8+0.25});
  }
  \fill[warncolor] (0,0) circle (3pt);
  \node[font=\small, below] at (0,-0.1) {\lr{$a$}};
  \node[font=\small, anchor=west] at (5.2,0.6) {\lr{$n=1$}};
  \node[font=\small, anchor=west] at (0.6,2.6) {\lr{$n$ \rl{بزرگ}}};
\end{tikzpicture}

\figcaption[fig:shrink]{بازه‌های کوچک‌شونده دور $a$؛ اشتراک همهٔ آن‌ها فقط نقطهٔ $a$ است.}
\end{figblock}

\textbf{بازه‌های بسته.} به همین ترتیب $[a,b]=\bigcap_n\bigl(a-\tfrac1n,\,b+\tfrac1n\bigr)$؛ پس $[a,b]$ \emph{قطعاً} باید در $\F$ باشد. (در لبه‌های $\Omega$ بازه‌ها را با $\Omega$ اشتراک بگیرید.)

\textbf{اعداد گویا و گنگ.} گویاها اجتماعی شمارا از نقطه‌ها هستند و گنگ‌ها متمم آن؛ پس هر دو در $\F$ هستند. (این همان محاسبهٔ مثال ۳ است که حالا جایگاه دقیقش را پیدا کرد.)

\subsection{تعریف میدان بورل}

برای تعریف دقیق، به مفهوم «کوچک‌ترین $\sigma$-میدانی که شامل یک خانوادهٔ داده‌شده است» نیاز داریم.

\begin{lemma}[label=lem:intersection]{اشتراک $\sigma$-میدان‌ها}
اشتراک هر تعداد $\sigma$-میدان روی $\Omega$ هم یک $\sigma$-میدان است.
\end{lemma}

\begin{proof}
$\Omega$ در همهٔ آن‌هاست، پس در اشتراک است. اگر $A$ در اشتراک باشد، در هر کدام از آن‌هاست؛ پس $A^c$ هم در هر کدام و در نتیجه در اشتراک است. برای اجتماع شمارا هم به همین ترتیب.
\end{proof}

\begin{definition}[label=def:generated]{$\sigma$-میدان تولیدشده}
اگر $\mathcal C$ خانواده‌ای از زیرمجموعه‌های $\Omega$ باشد، $\sigma(\mathcal C)$ را اشتراکِ همهٔ $\sigma$-میدان‌هایی تعریف می‌کنیم که $\mathcal C$ را در بر دارند. طبق لم \ref{lem:intersection}، $\sigma(\mathcal C)$ \emph{کوچک‌ترین} $\sigma$-میدانِ شامل $\mathcal C$ است و به آن \term{$\sigma$-میدان تولیدشده}{generated $\sigma$-field} توسط $\mathcal C$ می‌گوییم. به زبان ساده: هر چه با متمم و اجتماع و اشتراک شمارا از اعضای $\mathcal C$ ساخته می‌شود.
\end{definition}

\begin{definition}[label=def:borel]{میدان بورل}
$\sigma$-میدانی را که بر پایهٔ بازه‌ها ساخته می‌شود \term{میدان بورل}{Borel $\sigma$-field} می‌نامیم و با $\mathcal B$ نشان می‌دهیم:
\[ \mathcal B(\R)=\sigma\bigl(\setc{(a,b)}{a<b}\bigr),\qquad \mathcal B\bigl([0.5,1.5]\bigr)=\sigma\bigl(\setc{(a,b)}{0.5\le a<b\le1.5}\bigr). \]
اعضای آن را \term{مجموعه‌های بورل}{Borel sets} می‌گوییم.
\end{definition}

\begin{proposition}[label=prop:borel-members]{چه مجموعه‌هایی بورل‌اند؟}
مجموعه‌های زیر همه در $\mathcal B(\R)$ هستند (برای هر $a<b$ و هر $x\in\R$):
\begin{enumerate}
  \item هر نقطهٔ منفرد $\{x\}$؛
  \item بازه‌های $[a,b]$، $(a,b]$، $[a,b)$، $(a,\infty)$ و $(-\infty,a]$؛
  \item هر مجموعهٔ شمارا، به‌ویژه $\N$، $\Z$ و $\Q$؛
  \item مجموعهٔ اعداد گنگ $\R-\Q$.
\end{enumerate}
\end{proposition}

\begin{proof}
\begin{enumerate}
  \item $\{x\}=\bigcap_{n}\bigl(x-\tfrac1n,\,x+\tfrac1n\bigr)$.
  \item $[a,b]=\bigcap_{n}\bigl(a-\tfrac1n,\,b+\tfrac1n\bigr)$، \ $(a,b]=\bigcap_{n}\bigl(a,\,b+\tfrac1n\bigr)$، \ $[a,b)=\bigcap_n\bigl(a-\tfrac1n,\,b\bigr)$، \ $(a,\infty)=\bigcup_{n}(a,\,a+n)$ و $(-\infty,a]=(a,\infty)^c$.
  \item هر مجموعهٔ شمارا اجتماعی شمارا از نقطه‌هاست.
  \item متمم $\Q$ است.
\end{enumerate}
\end{proof}

\begin{figblock}
\begin{tikzpicture}
  \draw[graytext, thick, rounded corners=8pt, fill=graytext!4] (-6.6,-2.3) rectangle (6.6,2.3);
  \node[font=\small, graytext, anchor=north east] at (6.45,2.2) {\rl{همهٔ زیرمجموعه‌ها ($2^{\Omega}$)}};
  \draw[excolor, thick, fill=excolor!10] (-0.9,-0.05) ellipse (5.3 and 1.9);
  \node[font=\small, excolor!70!black, anchor=north] at (-0.9,1.8) {\rl{میدان بورل $\mathcal B$}};
  \node[font=\small] at (-4.1,0.55) {\lr{$(a,b),\ [a,b]$}};
  \node[font=\small] at (-4.4,-0.45) {\lr{$\{x\}$}};
  \node[font=\small] at (-2.6,-1.05) {\rl{گویاها، گنگ‌ها}};
  \node[font=\small] at (1.4,0.75) {\rl{مجموعه‌های باز و بسته}};
  \node[font=\small] at (-1.2,0.0) {\rl{اجتماع شمارای بازه‌ها}};
  \node[font=\small] at (1.9,-0.95) {\rl{مجموعهٔ کانتور}};
  \fill[warncolor] (5.55,-1.55) circle (2.5pt);
  \node[font=\small, warncolor] at (5.4,-0.95) {\rl{مجموعهٔ ویتالی}};
\end{tikzpicture}

\figcaption[fig:borel]{میدان بورل تقریباً هر مجموعه‌ای را که در عمل با آن سروکار داریم در بر دارد، ولی مجموعهٔ ویتالی در آن نیست؛ و همین به ما اجازه می‌دهد روی آن احتمال تعریف کنیم.}
\end{figblock}

\begin{note}
\textbf{مجموعهٔ ویتالی بورل نیست}؛ پس با انتخاب $\F=\mathcal B$، دیگر لازم نیست به آن احتمال بدهیم و تناقض قضیهٔ ویتالی پیش نمی‌آید. در نظریهٔ اندازه ثابت می‌شود (قضیهٔ \emph{لبگ}) که روی $\mathcal B\bigl([0.5,1.5]\bigr)$ دقیقاً یک قانون احتمال وجود دارد که به هر بازه طولش را بدهد. پس مدل مقاومت به‌طور کامل و بدون تناقض ساخته می‌شود:
\[ \Bigl(\ [0.5,1.5]\ ,\ \mathcal B\bigl([0.5,1.5]\bigr)\ ,\ \P([a,b])=b-a\ \Bigr). \]
\end{note}



%% ------------------------------------------------------------
\section{جمع‌بندی: یک مدل احتمال یعنی $(\Omega,\F,\P)$}
%% ------------------------------------------------------------

برای ساختن یک مدل احتمال باید سه چیز را تعریف کرده باشیم: فضای نمونه $\Omega$، $\sigma$-میدانِ پیشامدها $\F$، و قانون احتمال $\P$ روی $\F$. سه مثال این جلسه در یک نگاه:

\begin{center}
\renewcommand{\arraystretch}{1.5}
\begin{tabular}{>{\semifont}l c c l}
\toprule
مثال & $\Omega$ & $\F$ & $\P$\\
\midrule
سکه & $\{H,T\}$ (متناهی) & $2^\Omega$ & $\P(\{H\})=p$، \ $\P(\{T\})=1-p$\\
سرور & $\{0,1,2,\dots\}$ (شمارا) & $2^\Omega$ & $\P(\{i\})=2^{-(i+1)}$\\
مقاومت & $[0.5,1.5]$ (ناشمارا) & $\mathcal B\bigl([0.5,1.5]\bigr)$ & $\P([a,b])=b-a$\\
\bottomrule
\end{tabular}
\end{center}

در جلسهٔ بعد از همین سه اصل شروع می‌کنیم و قواعد مهم احتمال (متمم، یکنوایی، شمول و عدم شمول، کران اجتماع و پیوستگی) را از آن‌ها \emph{اثبات} می‌کنیم.

%% ------------------------------------------------------------
\section*{جمع‌بندی جلسهٔ دوم}
\addcontentsline{toc}{section}{خلاصه و تمرین‌ها}
%% ------------------------------------------------------------

\begin{summary}
\begin{itemize}
  \item برآمدهای آزمایش در فضای نمونهٔ $\Omega$ هستند؛ $\Omega$ ممکن است \textbf{متناهی}، \textbf{شمارا} یا \textbf{ناشمارا} باشد.
  \item روی مجموعه‌ها سه نگاه داریم: \textbf{توپولوژی} (فاصله و شکل)، \textbf{جبر} (عملیات) و \textbf{نظریهٔ اندازه} (وزن‌دهی). احتمال یعنی \emph{وزن دادن} به پیشامدها.
  \item \textbf{قانون احتمال:} $\P(A)\ge0$، \ $\P(\Omega)=1$، \ و جمع‌پذیری برای پیشامدهای مجزا.
  \item \textbf{سکه:} قانون‌ها $p$ را تعیین نمی‌کنند؛ معنای $p$ در دیدگاه \textbf{فراوانی} «نسبت تکرار» و در دیدگاه \textbf{بیزی} «میزان باور» است.
  \item \textbf{سرور:} جمع‌پذیری برای دو پیشامد کافی نیست؛ قانون ۳ باید برای اجتماع \textbf{شمارا} برقرار باشد.
  \item \textbf{مقاومت:} در $\Omega=[0.5,1.5]$ با «وزن = طول»، هر نقطه وزن صفر دارد؛ وزن گویاها صفر و وزن گنگ‌ها یک است. جمع‌پذیری ناشمارا به تناقض می‌رسد.
  \item \textbf{ویتالی:} روی فضای ناشمارا نمی‌توان همیشه روی کل $2^\Omega$ احتمال تعریف کرد؛ پس به یک \textbf{$\sigma$-میدان} عقب‌نشینی می‌کنیم: شامل $\Omega$ و بسته نسبت به متمم و اجتماع شمارا.
  \item قانون احتمال روی $\sigma$-میدان تعریف می‌شود، ولی هر $\sigma$-میدانی برای هر قانونی مناسب نیست. برای بازه‌های پیوسته، \textbf{میدان بورل} $\mathcal B=\sigma(\text{\rl{بازه‌ها}})$ را به کار می‌بریم که نقطه‌ها، همهٔ بازه‌ها، گویاها و گنگ‌ها را دارد.
  \item \textbf{فضای احتمال} $=(\Omega,\F,\P)$ و سه شرط $\P$ همان \textbf{اصول کولموگروف} است.
\end{itemize}
\end{summary}

\subsection*{تمرین‌های پیشنهادی}

\begin{exercise}
سکه‌ای طوری ساخته شده که احتمال شیر دو برابر احتمال خط است. با قانون‌های احتمال، $\P(\{H\})$ و $\P(\{T\})$ را بیابید.
\answer{$\P(\{H\})=2/3$ و $\P(\{T\})=1/3$}
\end{exercise}

\begin{exercise}
تاسی طوری ساخته شده که احتمال هر وجه متناسب با عدد روی آن است: $\P(\{i\})=c\,i$. مقدار $c$ و احتمال «زوج آمدن» را بیابید.
\answer{$c=1/21$ و $\P(\fa{زوج})=12/21=4/7$}
\end{exercise}

\begin{exercise}
در مدل سرور (مثال \ref{ex:server})، احتمال پیشامدهای (الف) «دست‌کم ۳ درخواست» و (ب) «تعداد درخواست‌ها فرد است» را حساب کنید. در هر مورد مشخص کنید از کدام قانون استفاده کرده‌اید.
\answer{(الف) $1/8$ \quad (ب) $1/3$}
\end{exercise}

\begin{exercise}
در مدل مقاومت (مثال \ref{ex:resistor})، احتمال اینکه مقدار مقاومت (الف) کمتر از $0.7$ یا بیشتر از $1.2$ اُهم باشد؛ (ب) دقیقاً $0.8$ یا دقیقاً $1.2$ اُهم باشد؛ (ج) یک عدد گنگ بین $0.9$ و $1.1$ باشد، چقدر است؟
\answer{(الف) $0.5$ \quad (ب) $0$ \quad (ج) $0.2$}
\end{exercise}

\begin{exercise}
دو موقعیت واقعی مثال بزنید که در آن‌ها دیدگاه فراوانی برای تعیین احتمال کارایی ندارد ولی دیدگاه بیزی معنا دارد. در هر مورد توضیح دهید «باور اولیه» از کجا می‌آید.
\end{exercise}

\begin{exercise}
روی $\Omega=\{1,2,3,4\}$، کوچک‌ترین $\sigma$-میدانی را که شامل $\{1\}$ و $\{2\}$ است کامل بنویسید. چند عضو دارد؟
\answer{۸ عضو (همهٔ اجتماع‌های «آجرهای» $\{1\}$، $\{2\}$ و $\{3,4\}$)}
\end{exercise}

\begin{exercise}
آیا اجتماع دو $\sigma$-میدان همیشه $\sigma$-میدان است؟ (راهنمایی: روی $\Omega=\{1,2,3\}$ این دو $\sigma$-میدان را در نظر بگیرید:
\[ \F_1=\bigl\{\varnothing,\{1\},\{2,3\},\Omega\bigr\},\qquad \F_2=\bigl\{\varnothing,\{2\},\{1,3\},\Omega\bigr\}.) \]
\answer{خیر؛ $\{1\}\cup\{2\}=\{1,2\}$ در اجتماع آن‌ها نیست.}
\end{exercise}

\begin{exercise}
ثابت کنید مجموعه‌های $[a,b)$، $(a,\infty)$ و $\{1,\tfrac12,\tfrac13,\dots\}$ بورل‌اند.
\end{exercise}

\begin{exercise}\hard
نشان دهید اگر در تعریف $\sigma$-میدان به جای «اجتماع شمارا» فقط «اجتماع متناهی» بخواهیم، ممکن است پیشامدهای مهمی از دست بروند: روی $\Omega=\N$، خانوادهٔ مجموعه‌هایی که یا متناهی‌اند یا متممشان متناهی است را در نظر بگیرید. نشان دهید این خانواده نسبت به متمم و اجتماع متناهی بسته است، ولی پیشامد «عدد زوج است» در آن نیست.
\end{exercise}
